\section{Very-general jet estimates}\label{sec:jets}

This section derives two contrasting consequences of the assumed
counterexample: an upper bound for the vanishing of scalar sections,
and strong jet separation on a projective bundle over the product.
Both arguments use the exclusion of covering curves of bounded genus
and normalized degree.

We retain the counterexample in Equation~\eqref{ex:counterexample}
and the induction hypothesis of Assumption~\ref{mmp:lower-models}.
Thus \(X\) is smooth projective,
\(K_X\) is pseudo-effective, and \(\kappa(X,K_X)=-\infty\), whereas
good minimal models exist in smaller dimensions. We use the geometric
exclusions of Section~\ref{sec:exclusions}. Dimension one is impossible, since a
smooth curve with pseudo-effective canonical divisor has genus at least
one.  Hence \(n=\dim X\ge2\).

Throughout the section, fix the late terminal \(\Q\)-factorial model
\(Y\) from
Proposition~\ref{prop:bounded-curves}.  Write \(K=K_Y\) and \(A=A_Y\).
The subsequent scaling models \(Y_t\) are small modifications of \(Y\);
the transforms \(A_t\) are effective up to \(\Q\)-linear equivalence,
and \(K_t+tA_t\) is nef, big, and semiample.  Put
\[
 \mu_t=\vol(X,K_X+tA)^{1/n}.
\]
Here, in the volume on \(X\), \(A\) denotes the original ample divisor.
We have \(\mu_t\to0\).  Choose an integer \(m>0\) with \(mK\) Cartier.

\subsection{Normalization and two local tools}

We first normalize the small adjoint divisors so that their volumes
stay bounded away from zero.  Fix a small rational number \(\epsilon>0\).
An integer \(q>0\) will be chosen after \(\epsilon\) and before \(t\).
As rational \(t\downarrow0\),
choose positive integers \(r=r(t)\) with
\begin{equation}\label{jet:normalization}
 r\mu_t\longrightarrow\epsilon,\qquad
 L=r(K+tA),\qquad L_b=L+bmK\quad(0\le b\le q).
\end{equation}
Only rational weights \(b\) are used for model constructions;
volume functions in integrals are extended continuously to real weights.
The positive part of \(L_b\) is
the nef semiample class
\[
 N_b=(r+bm)(K_s+sA_s)
 \quad\hbox{on }Y_s,\qquad s=\frac{rt}{r+bm}.
\]
When emphasizing the weight, denote this axis model by \(Y(b)=Y_s\)
and write \(A_b\) for its transform of \(A\).
All references to positive parts below mean these classes on their
scaling models, or their pullbacks to common resolutions.  Monotonicity
of volume in pseudo-effective order gives
\begin{equation}\label{jet:volume-comparison}
 \vol(L)\le \vol(L_b)\le
 \left(1+\frac{qm}{r}\right)^n\vol(L).
\end{equation}
Indeed \(L_b-L=bmK\) is pseudo-effective, and
\((1+bm/r)L-L_b=bmtA\) is effective up to equivalence.
Consequently \(N_b^n\) is bounded above and bounded away from zero,
uniformly for \(b\in[0,q]\), and tends uniformly to \(\epsilon^n\).
More explicitly, for every fixed \(q\), once \(t\) is sufficiently
small depending on \(q\),
\begin{equation}\label{jet:uniform-normalized-volume}
 \frac{\epsilon^n}{2}\le N_b^n\le2\epsilon^n
 \qquad(0\le b\le q).
\end{equation}
The constants in this display are independent of \(q\); only the
threshold for \(t\) depends on the fixed choice of \(q\).
All unspecified constants in this section may depend on
\(n,\epsilon,q,Y,A,m\), but not on sufficiently small \(t\) or on \(b\).

This normalization gives a convenient form of the curve exclusion.
A family of curves with uniformly bounded geometric genus and
\(N_b\)-degree cannot cover the corresponding \(Y_s\) for arbitrarily
small \(t\).  In fact \(s/t\to1\) uniformly in \(b\), and
\[
 (r+bm)\mu_s=\vol(L_b)^{1/n}
\]
is bounded above and away from zero.  Thus such curves would have
uniformly bounded \((K_s+sA_s)\)-degree divided by \(\mu_s\), contrary
to Proposition~\ref{prop:bounded-curves}.  We call this consequence the
\emph{normalized curve exclusion}.

We will apply the local estimates of Section~\ref{sec:common-tools}
on several different models. In each application, let \(Z\) denote the
ambient model of dimension \(d\), let \(P\) be its nef, big, semiample
class, and let \(V\) be a moving base component of dimension \(f\).
Lemma~\ref{common:moving-kernel} applies to the \emph{entire} kernels
of jets in one Cartier degree \(kP\), at levels separated by a gap
\(\delta\). It gives vanishing along \(V\) in the ordinary ideal power
of order
\[
 h=\lceil k\delta/2\rceil,\qquad k\delta\ge4,
\]
and, when \(Z\) is \(\Q\)-factorial and klt, the estimates
\begin{equation}\label{jet:tracking-estimates}
 P^f\cdot V\le (2/\delta)^{d-f}P^d,
 \qquad K_{V^*}P^{f-1}\le(K_Z+cP)|_V P^{f-1},
 \quad 0<c\le 2(d-f)/\delta\le2d/\delta.
\end{equation}
Here \(V^*\) is a smooth resolution of \(V\). The canonical estimate
uses Lemma~\ref{common:numerical-subadjunction}, which allows the nef
class on the center itself and requires log canonicity only near its
generic point. The last, weaker coefficient is sufficient here. At every
application we choose the incidence point in the isomorphism open of the
big semiample contraction. Dominance of the incidence permits this choice;
it ensures that $P|_V$ is big and hence $P^f\cdot V>0$.
Exceptional components not meeting that open are not assigned this
positivity. The separate fiber-restriction Lemma~\ref{common:fiber-restriction}
will be used only at a general incidence point where the scheme fiber is
reduced and the selected base component is isolated.

To apply normalized curve exclusion, we need curves with bounded
genus as well as bounded degree.  The preceding intersection estimates
provide them on components of general type.  Suppose a moving
\(f\)-fold \(V\) is of general type, \(P|_V\) is nef and big, and
\begin{equation}\label{jet:curve-input}
 0<P^f.V\le C_0,\qquad
 K_{V^*}P^{f-1}\le C_1P^f.V.
\end{equation}
Effective birationality gives a uniform pluricanonical degree whose
moving part, on a further resolution, is a big basepoint-free
Cartier divisor \(H\) defining a birational morphism
\cite[Theorem~4.0.1(3)]{HMX2018}.  Thus
\(HP^{f-1}\le C_2P^f.V\).  Mixed Hodge index gives
\[
 H^jP^{f-j}\le C_2^jP^f.V\quad(0\le j\le f);
\]
no lower bound on \(P^f.V\) is needed here.  For \(f>1\), cut by
\(f-1\) general members of \(|H|\). Their \(P\)-degree is
\(H^{f-1}P\le C_2^{f-1}C_0\), and their images under the birational
morphism defined by \(H\) have degree \(H^f\le C_2^fC_0\).
A general plane projection is birational on such an image curve;
the arithmetic genus of the plane image therefore bounds its geometric
genus in terms of this degree. For \(f=1\), the
canonical-degree bound already bounds the genus.
We will also use this construction for a log smooth pair of log
general type with reduced boundary, using coefficients in the fixed set
\(\{1\}\) in effective birationality.

If such components sweep an ambient space mapping to \(Y\), and
their curves project nonconstantly with \(N_b\)-degree bounded by
their \(P\)-degree, we obtain forbidden covering curves on \(Y_s\).
When the component has positive-dimensional image in the base, general
complete intersections from the big birational system project
nonconstantly. The components sweep, and their chosen incidence points
can lie outside any prescribed countable union of proper closed
subsets. Hilbert schemes and spaces of maps have only countably many
components, so these curves can be parameterized in covering algebraic
families. Thus, whenever Equation~\eqref{jet:curve-input} is
established below, it remains to check general type and the stated
covering and projection conditions to invoke normalized curve exclusion.

We first apply the two local tools to scalar sections.  Excessive
vanishing produces a proper moving component, whose general type
turns the numerical estimates into the forbidden curves.

\begin{proposition}[Scalar jet bound]\label{prop:scalar-jets}
For sufficiently small \(t\), let \(P\) be the positive part of
\(2r(K+2tA)\), and put \(\rho=(P^n)^{1/n}\).
At a very general point, every nonzero section of every Cartier
multiple \(kP\) has order at most \(4k\rho\).
The same statement holds on a common resolution after adding an
effective exceptional divisor that does not change the section
spaces.  Moreover
\[
 2r\mu_t\le\rho\le4r\mu_t,
\]
so in particular \(\rho\le8\epsilon\) for small \(t\).
\end{proposition}

\begin{proof}
First, pseudo-effectivity of \(K\) gives the volume bounds:
\[
 \vol(K+tA)\le\vol(K+2tA)\le2^n\vol(K+tA).
\]
Suppose the asserted order bound fails for arbitrarily small \(t\).
There are countably many Cartier degrees, and nonzero jet kernels in
each degree are detected by algebraic rank loci. Thus failure at a very
general point gives a degree in which a violating section exists on a
common constant-rank marked-point open. Taking powers makes this degree
\(k\) arbitrarily large and divisible. Choose \(n+1\) levels strictly
between \(\rho\) and \(4\rho\), with equal gaps comparable to \(\rho\).
The full jet kernels at these levels are nonzero. Write \(\tau_0\)
for the lowest level. If the marked point were isolated in its base
locus, local Bezout for \(n\) general kernel sections would contribute
at least \((k\tau_0)^n>k^nP^n\), exceeding the total intersection.
Lemma~\ref{common:moving-kernel} therefore supplies a moving proper
positive-dimensional component \(V\).

We can now apply the curve construction.  The component is of general
type by
Proposition~\ref{prop:general-type}. Here let \(Z\) be the scaling
model for \(K+2tA\), so that \(P=2r(K_Z+2tA_Z)\). The effective
transform \(A_Z\) does not contain a general such component, and
\(K_Z|_V\le(2r)^{-1}P|_V\) in effective order.
Equations~\eqref{jet:tracking-estimates} consequently imply
Equation~\eqref{jet:curve-input} with constants uniform in \(t\);
the gap is bounded above and away from zero because
\(\rho\) is comparable to the fixed \(\epsilon\).
The resulting bounded-genus, bounded-normalized-degree covering
curves contradict Proposition~\ref{prop:bounded-curves}.
The argument applies to every Cartier multiple, not only the divisible
degrees used for tracking: a section violating the bound in any Cartier
degree could be raised to a power in an arbitrarily divisible degree.
Its order and its degree increase by the same factor.
Finally, exceptional additions preserve the section spaces and the
orders at general points.  This proves the assertion on resolutions.
\end{proof}

\subsection{The two-slot bundle}

We next seek large jet separation by placing two copies of the
normalized adjoint divisor on a projective bundle.  The weight
decomposition of its sections will let us compare moving components
with either copy of the base.  On \(Y\times Y\), let
\begin{equation}\label{jet:two-slot}
 Z_0=\PP(mK_1\oplus mK_2),\qquad
 \xi=\OO_{Z_0}(1),\qquad M=L_1+L_2+q\xi.
\end{equation}
We use the convention that sections of \(k\xi\) are symmetric powers
of the indicated direct sum.  Fixing one base slot gives the
one-slot bundle
\(\PP(\OO_Y\oplus mK)\), with class \(L+q\xi\), up to a constant
line from the fixed slot.  We call it a \emph{slice}.
The torus open in either bundle is the complement of its two axes.

\begin{lemma}[Models, volume, and weights]\label{jet:bundle-models}
Both the total class \(M\) and the slice class have big semiample
positive parts \(P\) on \(\Q\)-factorial klt models \(Z\), with
\[
 K_Z+\Delta\sim_{\Q}c_tP,\qquad \Delta\ge0,
\]
where \(c_t\) is bounded independently of small \(t\).
Their volumes satisfy, respectively,
\begin{align}
 \vol(M)&=\frac{(2n+1)!}{(n!)^2}
   \int_0^q\vol(L_b)\vol(L_{q-b})\,db,\label{jet:total-volume}\\
 \vol(L+q\xi)&=(n+1)\int_0^q\vol(L_b)\,db.\label{jet:slice-volume}
\end{align}
In particular both volumes are bounded above and below by positive
constants times \(q\).  For each rational weight, on common
resolutions,
\begin{equation}\label{jet:weight-domination}
 P\sim_{\Q}N_{b,1}+N_{q-b,2}+E_b,\qquad E_b\ge0,
\end{equation}
with the analogous one-term formula on slices.
At a generic point of a base divisor of a slice, and on a general
torus fiber, the full system has no fixed subtraction.
\end{lemma}

\begin{proof}
We construct the semiample models as models of big klt adjoints.
The weight decomposition then gives both the volume formulas and
the comparison with the axis models.

The product \(Y\times Y\) is \(\Q\)-factorial. To see this, take a
smooth resolution \(\widetilde Y\to Y\). Its irregularity is zero by
Lemma~\ref{ex:irregularity}, so every line bundle on
\(\widetilde Y\times\widetilde Y\) is a tensor product of line bundles
from the two factors. The strict transform of a Weil divisor on
\(Y\times Y\) is therefore linearly equivalent to a sum of two
factorwise pullbacks. Push this relation down to \(Y\times Y\).
The resulting factor divisors are \(\Q\)-Cartier because \(Y\) is
\(\Q\)-factorial, and hence so is the original divisor. Products of
canonical singularities are canonical, and the projective bundles
are klt.  The two disjoint Cartier axes form a plt boundary \(D\),
and
\[
 K_{Z_0}+D=K_{\rm sum},\qquad
 D\sim2\xi-mK_{\rm sum}.
\]
The notation \(K_{\rm sum}\) means \(K_1+K_2\), or \(K\) on a slice.

For \(0<\sigma\le1\), put
\[
 \gamma=2\sigma+\frac{q(1+\sigma m)}r,\qquad
 \Xi_t\sim_{\Q}\xi+
       \frac{(1+\sigma m)t}{\gamma}A_{\rm sum}.
\]
Both endpoint weight systems are big, so choose an effective rational
representative of \(\Xi_t\) containing neither axis, and denote the
representative by \(\Xi_t\) as well.
To make the adjoint construction uniform, we first obtain a threshold
bound independent of \(\sigma\).
Restrict to \(0<t\le\sigma/(1+m)\).  Since \(\gamma\ge2\sigma\),
the coefficient
\[
 a=\frac{(1+\sigma m)t}{\gamma}
\]
lies in \((0,1]\).  Thus the numerical classes
\(\xi+aA_{\rm sum}\), and their restrictions to either fixed axis,
range over bounded segments independent of \(\sigma,q,r,t\).
On one fixed smooth resolution the effective pullbacks of these chosen
divisors have uniformly bounded degree against a fixed very ample class.
That degree bounds their multiplicity at every point, including
the coefficients of exceptional components.  Rational denominators
do not enter this estimate.  The smooth lc threshold is at least
the reciprocal of maximal multiplicity.

To transfer this estimate to the fixed klt space, write the crepant
boundary on its resolution as an SNC divisor with coefficients below one.
The minimum of
one and the positive numbers one minus its positive coefficients
bounds its log discrepancies below by a fixed positive multiple
of smooth log discrepancies.  Therefore the preceding smooth
bound gives a common klt threshold \(\eta>0\) for \(\Xi_t\) on the
original space. The same reasoning on each axis gives, after decreasing
\(\eta\), a klt threshold bound for the restricted divisor there.
The first bound will be used away from the axes; the restricted bounds
allow inversion of adjunction along the coefficient-one axes.

Now choose rational \(0<\sigma<\min\{1,\eta/4\}\), independently of
\(q,t\), and then take \(t\) small enough that
\(t\le\sigma/(1+m)\) and \(q(1+\sigma m)/r<\sigma\).
It follows that \(\gamma<3\sigma<\eta\).
Inversion of adjunction with the disjoint coefficient-one axes
gives plt near them; away from them the threshold bound gives klt.
Lowering their coefficients to \(1-\sigma\) consequently shows that
\[
 (Z_0,(1-\sigma)D+\gamma\Xi_t)
\]
is klt.  This choice respects the order: first \(\epsilon\), then
\(q\), then sufficiently small \(t\) with \(r\) as in
Equation~\eqref{jet:normalization}.
A direct calculation gives its adjoint class
\[
 K_{Z_0}+(1-\sigma)D+\gamma\Xi_t
   \sim_{\Q}\frac{1+\sigma m}{r}M.
\]
The big klt adjoint has a good minimal model by
\cite{BCHM2010}.  Pushing the boundary to this model gives the stated
\(\Delta\) and \(c_t=(1+\sigma m)/r\).

We turn to the volume formulas.  The weight-\(l\) summand in degree
\(k\) has base classes \(kL_{l/k}\) and \(kL_{q-l/k}\); on a slice there is one
such factor.  The section decomposition and asymptotic Riemann--Roch
give Equations~\eqref{jet:total-volume}--\eqref{jet:slice-volume}
as Riemann sums.  One can justify passage to the integral without
assuming uniform asymptotic Riemann--Roch: partition \([0,q]\) into
rational bins, compare weight classes in each bin by adding and
subtracting the bin width times a fixed very ample divisor
dominating both \(mK\) and \(-mK\), take divisible-degree limits,
and then shrink the bins.  Volume continuity gives the formulas.
Equation~\eqref{jet:volume-comparison} supplies their uniform bounds
and, in particular, proves bigness used above.
In the range of Equation~\eqref{jet:uniform-normalized-volume},
the bounds needed when choosing \(q\) are explicitly
\[
 \vol(M)\ge \frac{(2n+1)!}{4(n!)^2}\,q\epsilon^{2n},
 \qquad
 \vol(L+q\xi)\le2(n+1)q\epsilon^n.
\]
Their coefficients are independent of \(q\); the smallness threshold
for \(t\) is allowed to depend on \(q\).

It remains to establish the effective weight comparisons. On a common
resolution and in compatible divisible degrees, let \(F_{\mathrm{full}}\)
and \(F_b\) be the fixed divisors of the full system and its
weight-\(b\) subsystem, normalized by the degree. Since the second is a
subsystem, \(F_b\ge F_{\mathrm{full}}\). Its fixed divisor consists of
the toric monomial and the base-model fixed differences, and its moving
class is the sum of the indicated pullbacks of \(N_b\). Hence the full
moving class is that sum plus
\(E_b=F_b-F_{\mathrm{full}}\ge0\). This proves
Equation~\eqref{jet:weight-domination}.  The same comparison in
individual sufficiently divisible degrees shows that every section
at that weight, after the full fixed subtraction, contains the
corresponding multiple of \(E_b\).

At a generic base-divisor point the small base-model maps are
isomorphisms.  The two endpoint systems are free there in divisible
degree, so together they have no common zero on the projective-line
fiber.  The full system therefore has no fixed subtraction there.
The same holds on a general fiber.  Finally, when restricting
Equation~\eqref{jet:weight-domination} to a component sweeping the
torus open, choose a general component not contained in
\(\Supp E_b\).  There are only countably many rational weights, so
these effective restrictions can be required simultaneously.
\end{proof}

The bundle models are now available, and their volumes can be made
large by choosing \(q\).  We use this growth to prove a lower bound
for local positivity.  Recall that the Seshadri constant of a nef
\(\Q\)-Cartier class \(P\) on a projective variety \(Z\) at a smooth
point \(x\in Z\) is
\[
 \varepsilon(P;x)=
 \inf_{\substack{C\subset Z\ \mathrm{integral\ curve}\\x\in C}}
       \frac{P\cdot C}{\mult_x C}.
\]

\begin{proposition}[Large two-slot Seshadri constant]\label{prop:large-seshadri}
For each fixed sufficiently small \(\epsilon>0\), there is an integer
\(q>0\) such that, for all sufficiently small rational \(t>0\) and
\(r\) as in Equation~\eqref{jet:normalization}, the positive part
\(P\) of Equation~\eqref{jet:two-slot} has
\[
 \varepsilon(P;x)>2n+2
\]
at a very general smooth point of its torus open.
In particular, for some sufficiently divisible integer \(k>0\),
the complete system \(|kP|\) separates jets of order strictly
greater than \((2n+2)k\) at such a point.
The section spaces and these general-point jets agree on the
original bundle and on common birational resolutions.
\end{proposition}

\begin{proof}
The tracking lemma reduces failure of the assertion to a moving base
component.  We will first analyze its fibers over the two base factors,
and then treat components that are generically finite over both.

Write \(d=2n+1\). Choose a gap \(\delta>0\) large enough that
\[
 2(n+1)\epsilon^n(2/\delta)^n<1.
\]
This inequality will exclude a component that is a whole projective-line
fiber of a slice. Now choose \(q\) sufficiently large that
\[
 \left(\frac{(2n+1)!}{4(n!)^2}q\epsilon^{2n}\right)^{1/d}
 >2n+3+(d+1)\delta.
\]
For all sufficiently small \(t\), the \(d+1\) levels
\(\tau_i=2n+3+(i+1)\delta\), \(0\le i\le d\), then lie below
the volume root of the total positive part. Both choices precede
the choice of \(t\).
Assume that the Seshadri assertion fails for arbitrarily small \(t\).
Because the highest level is below the volume root, jet counts show
that the full kernels at all these levels are nonzero in sufficiently
large divisible degree. A curve through the marked
point with degree-to-multiplicity ratio below the lowest level is
contained in that system's base locus.  Such a curve exists from
the assumed Seshadri bound and the choice \(\tau_0>2n+3\), whether or
not the infimum defining the constant is attained. Thus the lowest
kernel has a positive-dimensional base component, and
Lemma~\ref{common:moving-kernel} supplies a
moving proper component \(V\) and both estimates
\eqref{jet:tracking-estimates}.

The numerical estimates already rule out components of general type.
For all moving components under consideration, the effective
boundary \(\Delta\) of Lemma~\ref{jet:bundle-models} does not contain
the component.  Thus \(K_Z|_V\le c_tP|_V\) in effective order.
Whenever \(V\) is of general type, the curve construction following
Equation~\eqref{jet:curve-input} and weight domination give a
contradiction.  To treat the remaining components, we begin with
their fibers over a base factor.

\paragraph{Restriction to a slice.}
Resolve the rational projection from the total model to the first copy
of \(Y\), and work on the open where this resolution and the original
bundle agree. Fixing the first coordinate there gives the opposite
slice. Write \(Z_{\mathrm{sl}}\) for its model from
Lemma~\ref{jet:bundle-models} and \(P_{\mathrm{sl}}\) for its positive
part; their dimension and volume are \(n+1\) and
Equation~\eqref{jet:slice-volume}, respectively. We use these slice
objects until returning to the total model below. Choose one sufficiently
divisible Cartier degree for the total model, the slice model, and
their comparison on a common resolution.
Choose a general incidence point on the common isomorphic torus open,
away from the other base components and where the ordinary-power
containment of Lemma~\ref{common:moving-kernel} holds. Generic smoothness
allows us also to require that \(V\) be smooth over its actual image
in the first slot there. Let \(F\) be the reduced component of the
scheme fiber through that point, and suppose \(0<\dim F<n+1\).
Lemma~\ref{common:fiber-restriction} now shows that, near that point, the
restricted higher-series base ideal has support exactly \(F\) and lies in the
ordinary power \(\mathcal I_F^h\), with \(h=\lceil k\delta/2\rceil\).
Because \(F\) is proper in the slice, the restricted series is nonzero
and its generic base ideal is primary.

The fixed differences of the full and slice models are units on a
further common open. After removing these fixed divisors and trivializing
the fixed-slot lines, all restricted higher-kernel sections form a
subspace of \(H^0(Z_{\mathrm{sl}},kP_{\mathrm{sl}})\), with the same
local base ideal. This subseries need not be the full jet kernel on
the slice: the numerical estimates come from
Lemma~\ref{common:base-component}, applied to the actual restricted
subseries.
The model \(Z_{\mathrm{sl}}\) is projective,
\(\Q\)-factorial and klt, and \(P_{\mathrm{sl}}\) is nef, big and
semiample. Choose the incidence point also in its smooth open.
The component \(F\) is proper, its generic base ideal is primary, and
the preceding restriction gives ordinary multiplicity at least
\(h\ge k\delta/2\). Writing \(f=\dim F\), the lemma therefore gives
\[
 \begin{aligned}
 P_{\mathrm{sl}}^f\cdot F
 &\le(2/\delta)^{n+1-f}P_{\mathrm{sl}}^{n+1},\\
 K_{F^*}P_{\mathrm{sl}}^{f-1}
 &\le(K_{Z_{\mathrm{sl}}}+cP_{\mathrm{sl}})|_F
       P_{\mathrm{sl}}^{f-1},
 \qquad 0<c\le2(n+1-f)/\delta.
 \end{aligned}
\]

The restricted components must also move sufficiently to apply curve
exclusion.  These \(F\)'s may be chosen in families sweeping the opposite
slice.  Indeed the incidence of the \(V\)-family dominates the
total torus.  Choose a general incidence point, not necessarily
the original marked point, and then a general fiber of its first
slot evaluation.  Dominance gives the required dominant evaluation
onto that fixed opposite slice. Choose a component with dominant
evaluation there; generic flatness permits the fiber-component choices
just made. For each fixed \(q,t\), very general choices
also avoid the exceptional sets for all rational weights.
Choose this point also in the isomorphism open of the slice's big
semiample contraction and outside the relevant fixed differences.
The restriction of its positive part to \(F\) is then nef and big,
so its top intersection on \(F\) is positive, as required for the
curve construction.

\paragraph{Proper base images on a slice.}
If \(F\) is generically finite over a proper positive-dimensional
image \(W\) in the remaining base, then \(W\), and hence \(F\), is
of general type by Proposition~\ref{prop:general-type}.
The slice estimates and subadjunction give the required bounds on
this component, while weight domination bounds the degrees of its
projected curves.  We therefore obtain forbidden bounded curves.

Over a proper base image, the other possibility is that \(F\) is
saturated over its base image \(W\): on the original bundle it is
the full projective-line bundle over \(W\).
Put \(w=\dim W=\dim F-1\).  The nef weight comparison gives
\begin{equation}\label{jet:saturated-degree}
 qN_b^w.W\le P_{\mathrm{sl}}^{w+1}.F.
\end{equation}
Indeed \(P_{\mathrm{sl}}-\pi^*N_b\) is effective on the resolved graph
of the slice projection \(\pi\) to the base, so
successive mixed nef intersections give
\(P_{\mathrm{sl}}^{w+1}.F\ge
P_{\mathrm{sl}}(\pi^*N_b)^w.F\). The latter is
\(qN_b^w.W\), since endpoint freeness gives fiber degree \(q\).
If \(W\) is a point, the lower bound \(q\) contradicts the slice
Bezout upper bound \(2(n+1)q\epsilon^n(2/\delta)^n\), because our
choice of \(\delta\) makes the latter strictly less than \(q\).

Suppose \(w>0\).  The ruled component \(F\) itself need not be of
general type, so we seek the adjunction estimate on its base \(W\).
We obtain it from the coefficients of the slice sections. Use the
section-space identification to view the subseries on the original
slice, where the toric weight decomposition is defined. The full fixed
factors are units at the chosen generic torus point, so this identification
preserves the order \(h\). In degree \(k\), expand each of these sections
into its toric weights. At the generic point of \(W\), let \(I_l\) be
the ideal generated by the coefficients of the nonzero weight \(l\)
over all these sections.  Put \(h=\lceil k\delta/2\rceil\).
All \(I_l\) lie in \(\mathfrak m_W^h\): vanishing along the generic
torus fiber says that the coefficient of every Laurent monomial
vanishes to this order.  Their sum is
\(\mathfrak m_W\)-primary.  Otherwise its larger common zero germ,
times the generic torus fiber, would contradict isolation of \(F\).

Our aim is to construct an effective boundary
\(\Theta\sim_{\Q}c'L_b\), for some rational weight \(b\), such that
\(W\) is an lc center near its generic point and
\(c'\le2(n-w)/\delta\). Numerical subadjunction will then give a
canonical-intersection bound on \(W\), even though \(F\) is ruled.
Work in the regular local ring at the generic point of \(W\), of
dimension \(c=n-w\). Resolve the finitely many coefficient ideals
simultaneously. The exponents for which their product defines an lc
ideal pair form the rational polytope
\[
 u_l\ge0,\qquad
 \sum_lu_l\ord_E(I_l)\le a(E;Y,0).
\]
Here \(E\) runs through the divisors needed to test log canonicity on
the simultaneous resolution, and \(a(E;Y,0)\) is the log discrepancy.
Include the exceptional divisor of the blowup of the closed point.
Since every \(I_l\subset\mathfrak m_W^h\), it gives
\(\sum_lu_l\le c/h\). The polytope is compact, and \(\sum_lu_l\)
has a positive rational maximum on it. Fix a rational maximizing point.
For each coordinate \(l\), some equality at this point must involve
a valuation with \(\ord_E(I_l)>0\); otherwise \(u_l\) could be
increased while staying in the polytope. Such a valuation has log
discrepancy zero for the ideal pair. Its center \(C_l\) is contained
in \(V(I_l)\) and contains the generic point of \(W\). Their intersection
is contained in \(V(\sum_l I_l)\), so it is exactly \(W\) locally.
To apply the intersection property for lc centers, we first realize
these ideal lc places by an actual divisor. Choose an integer
\(M_0>\max_l u_l\) and \(M_0\) general coefficient divisors
\(D_{l,j}\) from each system.  Set
\(\Theta=\sum_l(u_l/M_0)\sum_{j=1}^{M_0}D_{l,j}\).
On the simultaneous resolution its fixed crepant coefficients are
at most one, and its general moving divisors meet transversely
with coefficients below one.  Thus the pair is lc near the generic
point of \(W\), and every active valuation remains an lc place.
Each \(C_l\) is consequently an lc center of this pair. The intersection
property \cite[Theorem~1.7]{KollarKovacs2010}, applied to the identity
morphism on its lc open, shows that \(W\) is an lc center generically.
This also applies when some
maximizing coordinate \(u_l\) is zero: the active valuation blocking that
coordinate still has center in \(V(I_l)\).

Set
\[
 b=\frac{\sum_lu_l(l/k)}{\sum_lu_l},
 \qquad c'=k\sum_lu_l\le\frac{kc}{h}.
\]
Since a weight-\(l\) coefficient divisor has class \(kL_{l/k}\),
the constructed boundary has class \(c'L_b\). Also
\(c'\le kc/h\le2(n-w)/\delta\). Transform it and \(W\) to the small
scaling model \(Y(b)\); our general incidence choice places the
generic point of \(W\) in their common isomorphism open. On this
model the boundary has class \(c'N_b\). Numerical subadjunction
and \(N_b=(r+bm)K_{Y(b)}+rtA_b\) give
\[
 K_{W^*}N_b^{w-1}
 \le\left(\frac1{r+bm}+c'\right)N_b^w.W.
\]
Indeed a general \(W\) is not contained in a fixed effective
representative of \(A_b\), so its contribution to this intersection
is nonnegative. The slice degree estimate and
Equation~\eqref{jet:saturated-degree} give, in turn,
\[
 0<N_b^w.W
 \le\frac{(2/\delta)^{n-w}}q P_{\mathrm{sl}}^{n+1}
 \le2(n+1)\epsilon^n(2/\delta)^{n-w}.
\]
These are precisely the two bounds required in
Equation~\eqref{jet:curve-input}, with constants independent of \(t\).
The variety \(W\) is of general type by
Proposition~\ref{prop:general-type}; normalized curve exclusion
again gives a contradiction.

\paragraph{A slice multisection.}
Proper base images have now been excluded.  The remaining proper
positive-dimensional slice case is
\(\dim F=n\), with \(F\) generically finite of degree \(e\) over
the whole base.  Weight comparison and the slice degree bound give
\[
 eN_b^n\le P_{\mathrm{sl}}^n.F,
\]
so \(e\) is bounded uniformly in \(t\).
Here the two axes provide the boundary needed to use log general
type on the base.  The closure of \(F\) on the original bundle is
neither axis.
Intersect this closure with each Cartier axis and push the resulting
effective integral cycle to \(Y\); call the resulting Weil divisors
\(D_1,D_2\). The two axis classes differ by the pullback of \(mK\),
and \(F\) has generic degree \(e\) over \(Y\). The projection formula
therefore gives
\begin{equation}\label{jet:axis-signed}
 D_1-D_2\sim_{\Q} \pm emK.
\end{equation}
At the appropriate endpoint weight \(b\in\{0,q\}\), the difference
\(E_b|_F\) contains the corresponding toric-axis intersection with
coefficient \(q\).  There is no full fixed subtraction above
generic base-divisor points by Lemma~\ref{jet:bundle-models}.
Intersecting with \(N_b^{n-1}\) and using mixed nef comparison yields
\[
 qN_b^{n-1}.D_i\le P_{\mathrm{sl}}^n.F.
\]
The other divisor has bounded degree for the same \(N_b\), because
of Equation~\eqref{jet:axis-signed} and
\[
 0\le K_{Y(b)}N_b^{n-1}\le\frac{N_b^n}{r+bm}.
\]
We now apply the logarithmic curve construction on the base \(Y(b)\),
not on the multisection \(F\). The signed divisor
\(\pm(D_1-D_2)/(em)\) represents \(K_{Y(b)}\) by
Equation~\eqref{jet:axis-signed}. On a log resolution of this base, add
the reduced strict support of \(D_1+D_2\) and all exceptional divisors
to the canonical divisor. Proposition~\ref{prop:signed-alternative}
gives a big adjoint for the reduced strict support of \(D_1-D_2\) and
all exceptional divisors. Adding the remaining effective reduced boundary
shows that the resulting log canonical divisor is big as well.
Its intersection with the pulled-back \(N_b^{n-1}\) is bounded:
exceptionals vanish under projection, and reduced support has
degree no larger than \(D_1+D_2\).
Log effective birationality and the curve construction therefore
contradict normalized curve exclusion on the base itself.

\paragraph{Reduction to a correspondence.}
We now return to the total model \(Z\) and its positive part \(P\).
The slice analysis has excluded fiber dimensions strictly between
zero and \(n+1\) over either slot.  A full \((n+1)\)-dimensional fiber over
one slot would mean that \(V\) contains the whole opposite slice
over its first image \(W\).  Its fiber dimension over the other
slot would then be \(\dim W+1\), strictly between zero and \(n+1\),
unless \(V\) were the full total space.  That is impossible.
Hence \(V\) projects generically finitely in both slots and
\(\dim V\le n\).  If either image is proper, general type and the
total-space estimates give the previous curve contradiction.
The same is true if \(V\) is of general type.
Thus only the case \(\dim V=n\), dominant with bounded degree over both
copies of \(Y\), remains.  We next show that its branch divisors
cannot move; this will reduce the family to finitely many covers.

\paragraph{Moving branch divisors.}
Consider the dominating algebraic family of these correspondences,
resolving maps after shrinking the parameter space.  Suppose a
divisorial branch component of one projection moves.  On a
resolution \(V^*\), choose a ramified divisor \(R\) above its generic
point.  There is a rational weight \(b\) for which the effective
difference \(E_b|_{V^*}\) does not contain \(R\).
Indeed take a fixed free divisible degree of \(P\); some section
does not vanish generically on \(R\).  Monomial-weight sections
span that degree, so one of them does not vanish there.
Only finitely many weights are tested in this fixed family; their
model comparisons can be resolved simultaneously.

Let \(J\) be the branch divisor on the corresponding small axis
model. Restrict the effective difference to \(R\); this is permitted
because \(R\) is not in its support. Mixed nef intersections and
projection to \(J\) bound its \(N_b\)-degree by \(P^{n-1}.R\).
The ramification formula over the terminal target bounds the latter
by \(K_{V^*}P^{n-1}\): its other ramification terms are nonnegative,
and the pullback of the target canonical class is pseudo-effective.
The total-space subadjunction estimate now gives
\begin{equation}\label{jet:branch-degree}
 N_b^{n-1}.J\le P^{n-1}.R
 \le K_{V^*}P^{n-1}\le C_q.
\end{equation}
The ramification coefficient of \(R\) is a positive integer, so is
at least one; the other ramification and exceptional terms are
nonnegative.  The last bound is the total-space subadjunction
estimate.

To turn this degree bound into bounded curves, we still need a
canonical intersection bound on \(J\).  The required adjunction
does not assume that \((Y(b),J)\) is globally lc.
Terminal \(Y(b)\) is smooth in codimension
two.  Thus normalization and conductor adjunction along \(J\)
give
\[
 K_{J^\nu}+C=(K_{Y(b)}+J)|_{J^\nu},\qquad C\ge0
\]
in codimension one.  On resolving \(J^\nu\), exceptional divisors
map to codimension at least two and pair trivially with the
pulled-back \(N_b^{n-2}\).  Consequently, with
\(d_J=N_b^{n-1}.J\),
\begin{align}
 K_{J^*}N_b^{n-2}
 &\le (K_{Y(b)}+J)J N_b^{n-2}\notag\\
 &\le \frac{d_J}{r+bm}+\frac{d_J^2}{N_b^n}
 \le C'_q d_J.\label{jet:branch-adjunction}
\end{align}
For the second line, a general moving \(J\) is not a component of
a fixed effective representative of \(A_b\), giving the first
term; mixed Hodge index gives the second.  The lower bound for
\(N_b^n\) and Equation~\eqref{jet:branch-degree} give the final
uniform constant.

A moving \(J\) sweeps very general base points and is of general
type by Proposition~\ref{prop:general-type}.
Equations~\eqref{jet:branch-degree}--\eqref{jet:branch-adjunction}
therefore give forbidden bounded curves.  Branch components can
be named after a finite parameter extension and shrinking.
It follows that, for all sufficiently small \(t\), all divisorial
branch components in the chosen family are fixed.

\paragraph{Fixed covers.}
We now fix $t$. The branch complement is chosen from the one finite-type
family of correspondences under consideration, before any cover is
classified. After a finite parameter extension and shrinking, resolve
the two evaluation maps simultaneously over a smooth parameter open
and follow the finitely many
irreducible relative ramification divisors. The argument just given
excludes each component whose divisorial images sweep the base: such a
component would supply the bounded-genus covering curves of
Proposition~\ref{prop:bounded-curves}. For every remaining component the
closure of its divisorial evaluation image is a fixed proper closed
subset of $Y$, and the intersection estimates already bound the degrees
of both projections. These are precisely the hypotheses of
Lemma~\ref{common:finite-covers} for the smooth projective family of
resolved correspondences. It gives a smooth dense open $Y^0$ and only
finitely many finite normal Stein covers in either slot, all etale over
$Y^0$. Bad parameter fibers are removed by shrinking the parameter space,
not by deleting their possibly dominant images in $Y$. Both $Y^0$ and
the finite cover lists may depend on this fixed $t$; no common complement
or list as $t\to0$ is used.

For a general member \(V_a^*\) of the resolved family, factor the two
projections through their finite normal Stein covers. The preceding
finiteness lets us fix these as \(T_1\to Y\) and \(T_2\to Y\) after
restricting to a family whose images still dominate \(Y\times Y\).
Each map \(\beta_{i,a}:V_a^*\to T_i\) is birational, so the two
projections determine a birational map
\[
 \phi_a=\beta_{2,a}\circ\beta_{1,a}^{-1}:T_1\dashrightarrow T_2.
\]
The graphs of these maps form an algebraic family after choosing a
Hilbert-scheme component and shrinking for flatness and birationality
of both projections. Such a component with dominant evaluation exists:
there are only countably many graph Hilbert schemes and finitely many
cover choices, whereas the original incidence dominates the base
product. Its graph incidence has image of dimension \(2n\) in
\(T_1\times T_2\), because this product is finite over \(Y\times Y\).
The product is integral of dimension \(2n\), so the graph incidence
dominates it. Proposition~\ref{ex:fixed-cover} excludes exactly this
family of birational maps between the two fixed covers.

Every possible moving component \(V\) has now led to a contradiction,
so the Seshadri bound holds.  To obtain the assertion about jets,
observe that at a very general point the big semiample
contraction is an isomorphism onto its image near that point.
The jet interpretation of the Seshadri constant for the ample
class downstairs gives the stated divisible jet-separating
multiple.  Complete systems agree under the model comparisons,
so this conclusion transfers to the original torus open and to
common resolutions.
\end{proof}
